\subsection{Projective Garbling}
\label{app:garbling}

For the scheme of Section~\ref{sec:garbling}, the labels and common
encoding material have polynomial length.
Correctness requires that, for every polynomial bound on circuit size,
there is a negligible function $\varepsilon$ such that every circuit $C$
within that bound satisfies
\begin{equation}
\label{eq:garbling-correctness}
  \Pr_{(\GC,e)\sample\Garble(1^\lambda,C)}
  \bigl[\forall x\in\bits^n:\;
    \Eval(\GC,\Encode(e,x))=C(x)\bigr]
    \ge 1-\varepsilon(\lambda).
\end{equation}
We strengthen the fixed-input correctness clause of Hemenway et al.\
\cite[full version, Definition~1]{HemenwayEtAl2016AdaptiveGarbling}
by requiring one sampled garbling to be correct for all inputs
simultaneously.  This covers inputs chosen after seeing $\GC$.
The perfect-correctness convention of Garg and Srinivasan~
\cite[full version, Definition~3.7]{GargSrinivasan2018AdaptiveGarbling}
also implies this requirement.

\subsubsection{Adaptive privacy}
Write $\mathcal S=(\SimGarble,\SimEncode)$ for the simulator.
The experiments $\GarbleReal_{\mathcal A}(1^\lambda)$ and
$\GarbleIdeal_{\mathcal A,\mathcal S}(1^\lambda)$ initially give
$1^\lambda$ to the efficient adversary $\mathcal A$ and offer these
commands:
\begin{description}
  \item[$\GarbleCircuit(C)$.]
    Store $C$.  In the real experiment, compute
    $(\GC,e)\sample\Garble(1^\lambda,C)$.
    In the ideal experiment, compute
    $(\GC,\SimState)\sample
    \SimGarble(1^\lambda,\Topo(C))$.
    Return $\GC$, retaining $e$ or $\SimState$ privately.
  \item[$\EncodeInput(x)$.]
    For $x\in\bits^n$, compute $X\gets\Encode(e,x)$ in the
    real experiment.  In the ideal experiment, compute
    $X\sample\SimEncode(\SimState,C(x))$.  Return $X$.
  \item[$\Finish(b)$.]
    Terminate and output $b\in\bits$.
\end{description}
There is at most one accepted $\GarbleCircuit$ command, followed
by at most one accepted $\EncodeInput$ command.  Malformed,
out-of-order, or repeated commands return $\bot$ without changing state.
The adversary may terminate early; termination without an output bit is
interpreted as $\Finish(0)$.

\begin{definition}[Adaptive garbling privacy]
\label{def:adaptive-garbling}
A garbling scheme has adaptive simulation privacy if there exists an
efficient simulator $\mathcal S$ such that, for every efficient adversary
$\mathcal A$, there is a negligible function $\varepsilon_{\mathcal A}$
satisfying
\[
  \left|
    \Pr[\GarbleReal_{\mathcal A}(1^\lambda)=1]
    -\Pr[\GarbleIdeal_{\mathcal A,\mathcal S}(1^\lambda)=1]
  \right|\leq\varepsilon_{\mathcal A}(\lambda).
\]
\end{definition}

Thus the simulator produces $\GC$ before learning the output, and later
produces $X$ from $C(x)$ alone, without receiving $x$ or the gate
functions.  The adversary chooses $x$ after seeing $\GC$, but chooses
all of $x$ before receiving $X$.  This is the coarse-grained adaptive
privacy notion $\PrvOne$ of Bellare, Hoang, and Rogaway~
\cite{BellareHoangRogaway2012Adaptive}; it does not permit several
input encodings for the same garbling.  Here adaptivity concerns input
selection, not corruption of a garbler.

\begin{theorem}[Hemenway et al.]
\label{thm:adaptive-garbling}
If one-way functions exist, then projective garbling schemes with
adaptive simulation privacy exist in the standard model.
\end{theorem}
Adaptive privacy follows from the projective construction and
\cite[full version, Corollary~2]{HemenwayEtAl2016AdaptiveGarbling}.
Its concrete instantiation also satisfies~\eqref{eq:garbling-correctness}.
To see this, consider all four input-label combinations at every gate.
For the inner encryption in Appendix~A of that paper, decryption under
a wrong key accepts exactly when the final $\lambda$ bits of two PRF
outputs coincide.  The acceptance condition does not involve the
plaintext, even when that plaintext depends on a wire label.  For independently sampled
labels its probability is negligible by PRF security.  A union bound
over the polynomially many row tests and wire-label collisions gives
one event outside which every gate evaluates correctly on every valid
label combination.  Induction over gates then gives correctness for
all inputs simultaneously.  The outer encryption has perfect
correctness~\cite[full version, Definition~2, Figure~1, and Appendix~A]
{HemenwayEtAl2016AdaptiveGarbling}.
